\section{Instance segmentation with a custom Mask R-CNN}
\label{sec:maskrcnn}

The first learned formulation treats lunar features as \emph{instances}: for every tile, a list of individual feature objects, each with its own class label, bounding box, and pixel-level mask. Mask R-CNN \citep{he2017maskrcnn} is the natural tool, because it predicts a segmentation mask in parallel with the detection of each object.

\paragraph{Instance-target preparation.}
A semantic channel is not the same thing as a set of objects. Before training, every channel is audited: only channels labelled in at least 200 tiles are kept, and channels whose mean coverage exceeds 20\% of the tile are flagged as coverage-like rather than object-like. Training therefore uses four classes: impact crater, pit/skylight, wrinkle ridge, and lobate scarp. Each selected channel is converted into instances by connected-component labelling (8-connectivity); components smaller than 30\,px or larger than 6000\,px are discarded, slivers narrower than 3\,px are dropped, and instances are capped at 100 per tile. This step is deliberately conservative, because noisy targets propagate directly into poor detections. A spatial band (the right-most tile columns) is held out for validation, and training tiles are augmented with geometric transforms plus brightness/contrast jitter.

\paragraph{Architecture and training.}
The backbone is a COCO-pretrained ResNet-50 with a Feature Pyramid Network \citep{he2016resnet,lin2017fpn}, lowest layers frozen and the top three trainable. Three changes adapt it to lunar tiles: anchors retuned to lunar scales (one size per FPN level, $(8,16,32,64,128)$\,px, aspect ratios $(0.5,1,2)$, preserving the pretrained RPN head shape); a deeper four-layer box head (1024 units, dropout 0.10); and a residual convolutional mask head (four residual blocks, 256 channels, GroupNorm) with transposed-convolution upsampling. The native 256\,px resolution is kept instead of torchvision's 800\,px rescaling. Optimisation uses AdamW \citep{loshchilov2019adamw} with discriminative learning rates (backbone and RPN at one tenth of the head rate of $2\times10^{-4}$), weight decay $10^{-4}$, one-epoch warm-up with cosine decay, and gradient clipping at 5.0, for 20 epochs at batch size 8. Because the only available hardware was a CPU (${\sim}$half a day per run), the loop was engineered to be crash-safe and resumable, checkpointing model, optimiser, and scheduler state every epoch.

\paragraph{Results.}
Table~\ref{tab:metrics} collects the validation metrics of the best checkpoint (epoch 8; validation loss turns upward after ${\sim}4$ epochs, indicating overfitting on the limited positives). The two halves of the network tell very different stories. The \emph{segmentation} branch works: whenever a detection is matched to a ground-truth object, the predicted mask overlaps it well (IoU 0.64, Dice 0.77). The \emph{detection} branch is where everything breaks down: mAP@0.5 sits near 1\%, and false negatives (7790) vastly outnumber true positives (461). The per-class breakdown shows that the comparatively common craters and ridges carry essentially all of the average precision, while the scientifically interesting pit/skylight class is barely detected and lobate scarps reach exactly zero --- too few examples to learn from. The training history and qualitative prediction panels are collected in Appendix~\ref{app:maskrcnn}.

\begin{table}[htbp]
\caption{Mask R-CNN validation metrics (best checkpoint, full held-out spatial band), with per-class AP@0.5.}
\label{tab:metrics}
\centering
\small
\begin{tabular}{l r}
\toprule
Metric & Value \\
\midrule
mAP@0.5                       & $0.010$ \\
Precision / Recall / F1       & $0.104$ / $0.056$ / $0.073$ \\
Mean mask IoU / Dice (matched) & $0.636$ / $0.773$ \\
Count MAE (objects per tile)  & $6.30$ \\
True / False pos.\ / False neg.\ & $461$ / $3973$ / $7790$ \\
\midrule
AP@0.5: impact crater          & $0.0233$ \\
\phantom{AP@0.5:} wrinkle ridge   & $0.0172$ \\
\phantom{AP@0.5:} pit / skylight  & $0.0003$ \\
\phantom{AP@0.5:} lobate scarp    & $0.0000$ \\
\bottomrule
\end{tabular}
\end{table}

\paragraph{Lessons.}
Most of what limits the score lies outside the network. Connected-component labelling is a leaky abstraction --- touching craters merge, elongated ridges fragment --- so the supervision is fundamentally noisier than a hand-annotated detection dataset, capping achievable precision. The rarest classes are the scientifically most interesting ones, and frequency-based channel selection cannot manufacture positive examples that are not there. Most instructively, good masks did not buy good detections: mask IoU is measured only on already-matched objects, so the bottleneck is upstream, in proposal recall and classification --- improving the mask head further would not move the headline metric.
